\section{Research questions and claim map}
\label{sec:questions}

The paper is organized around six questions. Each question has a primary
estimand, a diagnostic that could block the claim, and an evidence status. This
structure prevents the experiments from appearing as an unrelated plot catalog.

\begin{table}[tbp]
\centering\small
\begin{tabularx}{\linewidth}{@{}p{7mm}p{36mm}p{38mm}Yp{18mm}@{}}
\toprule
RQ & Question & Required diagnostic & Main artifact & Status \\
\midrule
1 & Does an order parameter retain semantic meaning? & latent/shuffled probe;
separate IID and OOD risk & observable dictionary and reference grid &
\statussupported{} \\
2 & Does a response sharpen into a finite-size boundary? & interior peak or
crossing; window and size sensitivity & GPU reference diagnostics &
\textcolor{AtlasRed}{\textbf{negative}} \\
3 & Which module or interaction causally supports task order? & matched ablation;
non-vanishing normalization; fresh seeds & MLP confirmation & \statusbounded{} \\
4 & How do optimizer and scale path change the learned state? & frozen selection
slice; compute and cap audit & optimizer and scaling confirmations &
\statusbounded{} \\
5 & Can a calibrated relation predict an untouched condition? & nontrivial
baselines; disjoint holdout; censoring & controlled continuation simulator &
model validation only \\
6 & Which conclusions survive a natural or learned endpoint? & comparable
estimand, disjoint split, independent implementation & corpus and cross-domain
dossiers & partial; mostly open \\
\bottomrule
\end{tabularx}
\caption{Question-to-claim map. ``Bounded'' means that the registered contrast
is supported while a neighboring stronger interpretation remains unestablished.}
\label{tab:rq-map}
\end{table}

The reporting template is the same for every result:
\emph{observation} states the estimate; \emph{interpretation} states the narrow
reading justified by its null and evidence vector; \emph{not established}
blocks a tempting stronger claim; and \emph{next falsifier} identifies the
experiment that could change the status. Sections~\ref{sec:results}--
\ref{sec:crossdomain} follow this order.

\FloatBarrier
